\section{Discussion and Conclusion}

\subsection{Summary of findings}

Our experiments tested whether the learned inverse-observation idea of
\citet{frerix2021invobs} survives a more operationally realistic setting: the
full four-dimensional variational data assimilation (4D-Var) cost
$J = \Jb + \Jo$ with an explicit background term and observation-error
covariance, Gaussian observation noise on every measurement, and an
operational sliding-window (cycling) extension. The central claim of the
paper broadly held. On Lorenz-96, under the tuned-$\sigb$ full 4D-Var cost,
learned-inverse-observation initialization combined with hybrid
(physics-space then observation-space) optimization attained the lowest
analysis root-mean-square error (RMSE) at every tested noise level
($\rmse = 0.620$, $1.168$, and $1.676$ at $\sigobs = 0.1$, $0.5$, and $1.0$
respectively). Equally important, hybrid optimization rescued the cheap
climatological-mean baseline from the bad observation-space local minima that
trap pure observation-space 4D-Var when only $10$ of $40$ components are
observed: it roughly halved the analysis error (for example $2.393 \to 1.223$
at $\sigobs = 0.5$). The same mechanism appeared in the sliding-window
experiments, where baseline initialization with observation-only optimization
remained stuck near the climatological RMSE while every hybrid-warm-started
method reached an analysis RMSE near $1$ at $\sigobs = 0$. Cycling delivered a
real but conditional benefit: the per-cycle analysis dropped below the
single-window reference once history accumulated, and the advantage was
clearest at higher noise and longer observation horizons
($\sigobs = 0.5$, $T_{\mathrm{obs}} = 48$), while being marginal at
$\sigobs = 0$, where a single well-posed window already over-determines the
observed components. On 2D Kolmogorov flow the qualitative story carried over:
learned-inverse-observation initialization placed L-BFGS in a basin roughly
$3\times$ lower in cost than bicubic interpolation (final losses $56{,}385$
and $45{,}760$ for the invobs combinations versus $168{,}341$ and $138{,}254$
for the baseline), and produced the best in-window analyses. In every case,
once the chaotic forecast left the assimilation window all methods relaxed to
the same climatological saturation, consistent with the short measured
Lyapunov time of the system---analysis quality improves the short-lead
forecast, not the Lyapunov-limited long-range predictability.

\subsection{Caveats and limitations}

This is a course capstone, and several limitations should temper how strongly
its conclusions are read.

\emph{Small ensembles.} The tuned-$\sigb$ Lorenz-96 sweep used only $N = 16$
trajectories, the sliding-window runs used $N = 50$, and the Kolmogorov
sliding-window runs used only $N = 8$. Consequently several ``best''
selections---the best stride, the best single-window method at $\sigobs = 0$,
and the cycling-versus-single-window margin in the clean-observation
cases---fall within ensemble scatter and should be read as tendencies rather
than decisive separations.

\emph{Compute-limited Kolmogorov runs.} The Kolmogorov solver relies on
\texttt{torch.fft.rfft2}, which ran on the CPU only (no graphics processing
unit, GPU), so all 2D experiments were validation-scale: a tiny ensemble, a
lightweight inverter trained on only a few hundred trajectories for a few tens
of epochs rather than the paper-scale $32{,}000$ trajectories and $500$
epochs, and a single observation-noise level. These runs demonstrate the
pipeline mechanics and reproduce the qualitative basin-of-attraction benefit,
but they are not converged scientific results, and the cycling machinery did
not yet separate from a single well-initialized window at this scale.
Relatedly, the warmup diagnostic reached only approximate quasi-stationarity
(a last-$20$-step enstrophy drift of about $10.8\%$, marginally above the
$\lesssim 10\%$ heuristic), a direct artifact of the short single-trajectory
diagnostic.

\emph{Inverter trained at one noise level.} The learned inverse observation
operator $\Hinv$ is supervised at a fixed noise level and then, in some cells,
evaluated under a deliberate mismatch. This probes robustness but means the
inverter is never optimal for the noise it actually encounters; a per-noise
checkpoint was supplied only where feasible.

\emph{Short measured Lyapunov time.} Our fixed-step fourth-order
Runge--Kutta (RK4) integrator yields a fitted Lyapunov time
$T_L = 0.442$ model time units (MTU), shorter than the paper's expected
$\sim 0.6$ MTU. We report this honestly as a known consequence of the coarse
fixed step being slightly too coarse to resolve the fastest unstable
directions. It sets the chaotic saturation rate to which every method relaxes
in the forecast and slightly tightens the window over which our RK4 and an
adaptive Dormand--Prince (dopri5) integrator represent the same physics
(median crossover $t^\star = 7.70$ MTU); our assimilation windows
($T = 1.0$ MTU) sit comfortably inside that horizon, so the integrator choice
does not contaminate the data assimilation (DA) results.

\emph{The negative $\sigb$-ordering result.} We expected
invobs$+$hybrid to prefer a \emph{smaller} optimal background weight $\sigb$
than the paper baseline, since the learned lift should make the analysis trust
the data more. This did not hold: invobs$+$hybrid selected
$\sigb = 10.0/1.0/1.0$ while paper$+$obs selected $\sigb = 0.1/0.3/0.3$, the
opposite direction at all three noise levels. Moreover, tuning $\sigb$ helped
only marginally over the untuned $\sigb = 1$ in most cells (gains typically
below $0.2$ RMSE, and exactly zero where the best $\sigb$ was already $1.0$).

\emph{No real data and simplified covariances.} All experiments are on
synthetic, perfect-model trajectories with no model error, and our covariances
are deliberately simple: a single static spatial-correlation matrix $\bm C$
(condition number $8.72$) for the Lorenz-96 background $\Bcov = \sigb^2 \bm C$,
a diagonal observation-error covariance $\Rcov = \sigobs^2 \bm I$ with
independent noise, and for Kolmogorov flow an identity background
$\Bcov = \sigb^2 \bm I$ because forming and factorizing a full
$4096 \times 4096$ spatial covariance was prohibitively expensive. None of
these capture the flow-dependent, correlated error structures of an
operational system.

\subsection{Future work}

Each limitation suggests a concrete next step. The inverter could be made
\emph{noise-aware}---either a family of per-noise checkpoints or a single
network conditioned on $\sigobs$---so that $\Hinv$ is always matched to the
observation error it encounters. The Kolmogorov pipeline should be moved to a
GPU and rerun at paper scale, with much larger ensembles, a fully trained
inverter, and a sweep over multiple noise levels, so that the cycling
extension can be evaluated where it has a chance to separate from a single
window. On the statistical side, both covariances invite enrichment: a
\emph{correlated} observation-error covariance $\Rcov$ in place of the diagonal
$\sigobs^2 \bm I$, and a flow-dependent background covariance estimated from an
ensemble (an ensemble or hybrid ensemble-variational scheme) in place of the
static $\bm C$, which is also the natural route to a tractable background
covariance for the 2D flow without forming the full dense matrix. Finally, an
adaptive integrator (dopri5) would produce paper-matched trajectories with the
expected Lyapunov time, removing the fixed-step artifact and enabling a cleaner
quantitative comparison against \citet{frerix2021invobs}; larger ensembles
throughout would turn the present tendencies into statistically firm claims.

\subsection{Conclusion}

Across both a 1D chaotic toy model and a 2D turbulent flow, and under a full
4D-Var cost with noisy observations that the original study did not consider,
the learned inverse observation operator continued to help where it was
designed to: it provides a physics-space initialization and a near-convex
hybrid objective that together place gradient-based 4D-Var in a good basin and
yield the lowest analysis error, especially when sparse, lossy observations
would otherwise trap pure observation-space optimization. The operational
sliding-window extension showed that cycling can further reduce analysis error
as history accumulates, most clearly when observations are noisy. At the same
time we report honest negatives---an unmet expectation about the optimal
background weight, a short integrator-limited Lyapunov time, and
compute-limited 2D runs---that mark exactly where larger ensembles, GPU-scale
Kolmogorov experiments, richer covariances, and an adaptive integrator would
turn these encouraging tendencies into firm conclusions.
